\uuid{7aqI}
\chapitre{Statistique}
\niveau{L2}
\module{Analyse de données}
\sousChapitre{Tests d'hypothèses, intervalle de confiance}
\titre{Expérience scientifique}
\theme{statistiques, tests d'hypothèses}
\auteur{}
\datecreate{2022-09-28}
\organisation{AMSCC}
\difficulte{}
\contenu{
\texte{Les quantiles sont à gauche : $q_{\nu;\gamma}$ pour la loi $\chi^2(\nu)$, avec une probabilité $\gamma$ à gauche du quantile.}
\texte{Une étude historique présente cinq catégories d'avance ou de retard scolaire, notées $-2,-1,0,+1,+2$. Pour cet exercice, on fixe par convention les probabilités de référence
\[(0{,}06;0{,}23;0{,}48;0{,}22;0{,}01).\]
Ces nombres sont inspirés des fréquences publiées par Binet sur 192 enfants ; on les traite ici comme une référence fixée, sans estimer son incertitude. Deux échantillons indépendants donnent :
\begin{center}\begin{tabular}{lrrrrr}\hline
Catégorie & $-2$ & $-1$ & 0 & $+1$ & $+2$ \\
Goddard (1547 enfants) & 294 & 309 & 557 & 325 & 62 \\
Bobertag (228 enfants) & 6 & 40 & 119 & 57 & 6 \\ \hline
\end{tabular}\end{center}}
\question{Tester séparément l'adéquation de chaque distribution à la référence au seuil de $5\%$. Vérifier les effectifs attendus et regrouper les classes si nécessaire.}
\indication{Les degrés de liberté sont le nombre de classes retenues moins un, car les probabilités de référence sont fixées.}
\reponse{Sous $H_0$, $e_j=np_j$ et $Q=\sum\frac{(O_j-e_j)^2}{e_j}$ est approximativement de loi $\chi^2(c-1)$.
Pour Goddard, les effectifs attendus sont $92{,}82;355{,}81;742{,}56;340{,}34;15{,}47$, tous supérieurs à 5. On obtient $q_{\mathrm{obs}}\simeq629{,}2123$, 4 ddl et $p_{\mathrm{val}}\simeq7{,}37\times10^{-135}$ : rejet à $5\%$.
Pour Bobertag, l'effectif attendu de $+2$ est seulement $2{,}28$. On regroupe $+1$ et $+2$. Les effectifs observés deviennent $(6,40,119,63)$ et les attendus $(13{,}68;52{,}44;109{,}44;52{,}44)$.
On trouve $q_{\mathrm{obs}}\simeq10{,}22424$, 3 ddl, $p_{\mathrm{val}}\simeq0{,}016753$ et $q_{3;0{,}95}\simeq7{,}815$ : rejet également.
Ces tests ne constituent pas une comparaison exacte à une population française estimée sur 192 enfants : pour cette autre question, il faudrait les effectifs français complets et un test d'homogénéité.}
}
